% Generated draft: requires mathematical and source review.
\documentclass[11pt]{article}
\usepackage[margin=25mm]{geometry}
\usepackage{fontspec}
\setmainfont{Latin Modern Roman}
\usepackage{amsmath,amssymb,mathtools,unicode-math}
\setmathfont{Latin Modern Math}
\usepackage{enumitem,needspace,xurl,hyperref}
\setlength{\parindent}{0pt}
\setlength{\parskip}{5pt}
\setlength{\emergencystretch}{3em}
\newcommand{\N}{\mathbb{N}}
\newcommand{\Z}{\mathbb{Z}}
\newcommand{\Q}{\mathbb{Q}}
\newcommand{\R}{\mathbb{R}}
\newcommand{\C}{\mathbb{C}}
\newcommand{\F}{\mathbb{F}}
\newcommand{\A}{\mathbb{A}}
\newcommand{\PP}{\mathbb{P}}
\newcommand{\E}{\mathbb{E}}
\newcommand{\eps}{\varepsilon}
\newcommand{\dd}{\,\mathrm{d}}
\begin{document}
\section*{UnsolvedMath ID: 800220}
\textbf{Source problem code:} catalogue entry
\subsection{Definitions and mathematical statement}
Let \(n\geq 1\) be an integer, and let \(\Omega\Subset\C^n\) be a domain with \(C^\infty\) boundary. Thus \(\Omega\) is connected and open, and its closure is compact. Assume that \(\Omega\) is pseudoconvex; for this smooth boundary, this means that the Levi form of a local smooth defining function, chosen negative on \(\Omega\), is nonnegative on the complex tangent space at each boundary point.

Write
\[
A^\infty(\bar\Omega)=\mathcal O(\Omega)\cap C^\infty(\bar\Omega),
\]
where \(\mathcal O(\Omega)\) denotes the holomorphic functions on \(\Omega\), and smoothness up to the boundary means that every real partial derivative extends continuously to \(\bar\Omega\). Say that \(\bar\Omega\) has a Stein neighborhood basis if every open set \(W\subset\C^n\) containing \(\bar\Omega\) contains a Stein open set \(V\) satisfying \(\bar\Omega\subset V\subset W\).

\textbf{Question.} If \(\bar\Omega\) has a Stein neighborhood basis, is it true that for every \(f\in A^\infty(\bar\Omega)\) and every \(\varepsilon>0\), there exist an open neighborhood \(U\subset\C^n\) of \(\bar\Omega\) and a function \(g\in\mathcal O(U)\) such that
\[
\sup_{z\in\bar\Omega}|f(z)-g(z)|<\varepsilon?
\]
The neighborhood \(U\) may depend on \(f\) and \(\varepsilon\); the requested approximation is uniform on \(\bar\Omega\).
\subsection{Short English statement}
For a bounded smooth pseudoconvex domain whose closure has a Stein neighborhood basis, can every holomorphic function smooth up to the boundary be approximated uniformly on the closure by functions holomorphic on neighborhoods of that closure?
\subsection{Sources}
\begin{itemize}
\item Local UnsolvedMath registry, permanent ID \(800220\), source problem code \texttt{catalogue entry}, titled \emph{Several Complex Variables Problem List --- Problem 16}. The supplied registry labels the problem open; that label is not independently verified here.
\item Buzzard, Fornaess, and Krantz, \emph{The Several Complex Variables Problem List} (1995), Problem 16. Supplied source URL: \url{https://arxiv.org/abs/math/9509202}. The record reports access on 2026-10-16 and extraction from arXiv LaTeX; these are supplied provenance claims, not a claim of access or verification in preparing this draft.
\item The source statement attributes this question to N. Sibony and supplies the contact address \texttt{sibony@anh.matups.fr}. The registry separately lists Buzzard, Fornaess, and Krantz in its proposer field, with year 1995.
\end{itemize}
\end{document}
